\documentclass[11pt]{article}
\usepackage[T1]{fontenc}
\usepackage{lmodern,amsmath,amssymb,amsthm,booktabs,longtable,microtype}
\usepackage[margin=1in]{geometry}
\usepackage[colorlinks=true,linkcolor=blue,urlcolor=blue,citecolor=blue]{hyperref}
\newtheorem{theorem}{Theorem}
\newtheorem{lemma}[theorem]{Lemma}
\newtheorem{proposition}[theorem]{Proposition}
\newtheorem{corollary}[theorem]{Corollary}
\newcommand{\Cyc}{\mathcal C}
\DeclareMathOperator{\comp}{comp}
\setlength{\emergencystretch}{3em}
\title{The exact cycle bound for loopless four-regular multigraphs}
\author{Research working draft}
\date{10 November 2026}
\begin{document}
\maketitle

\begin{abstract}
We prove that a loopless four-regular multigraph on $n$ vertices has at most $2^n+n$ cycles.
Equality holds precisely for doubled cycles and $K_5$.
We also classify all graphs above the one-half cycle-space bound and obtain an Eulerian rank bound.
The proof combines cycle-space reductions, a parity count, and a published degree-sequence path bound.
The remaining finite verification counts $925$ multigraph instances on at most ten vertices.
Two independent counting methods agree on these instances and on all tested instances through order eleven.
The finite verification is reproducible; this proof has not been formalised in Lean.
\end{abstract}

\section{Statement and conventions}

A cycle is a nonempty connected edge set whose incident vertices all have degree two.
Parallel edges are distinct.
Thus two parallel edges form a cycle of length two.
Write $c(G)$ for the number of cycles, and $P_{st}(G)$ for the number of vertex-simple paths between distinct vertices $s,t$.
The doubled $n$-cycle replaces each edge of a cycle by two parallel edges.
For $n=2$, this means the graph with two vertices and four parallel edges.

\begin{theorem}\label{thm:main}
Every finite loopless four-regular multigraph $G$ on $n$ vertices satisfies
\[
 c(G)\le2^n+n.
\]
Equality holds if and only if $G$ is a doubled cycle or is isomorphic to $K_5$.
\end{theorem}

The inequality is Open Problem~2 in Section~6 of Alahmadi, Aldred, dela Cruz, Sol\'e, and Thomassen~\cite{aacst}.
Their Theorem~6 already gives $c(G)\le2^n$ unless $G$ has a Hamilton decomposition.
We include the short cycle-space argument that gives this reduction.
Our analytic input is Theorem~7 of Dvo\v r\'ak, Morrison, Noel, Norin, and Postle~\cite{dmnnp}.
The finite computation uses \texttt{geng} from nauty~2.8.9~\cite{nauty}.
The proof below does not require contraction of an arbitrary doubled edge.

For a graph $F$, put $\beta(F)=|E(F)|-|V(F)|+\comp(F)$.
Its binary cycle space $\Cyc(F)$ consists of all even edge sets and has size $2^{\beta(F)}$.
Vertices of degree zero are permitted when computing this space.

\section{Reductions and equality under the reductions}

\begin{lemma}[Terminal paths]\label{lem:paths}
Let $N$ be connected and loopless on $s\ge2$ vertices.
Suppose that distinct terminals $u,v$ have degree two and all other vertices have degree four.
Then
\[
 P_{uv}(N)\le2^{s-1}.
\]
Equality holds only when $N$ is a doubled Hamilton path from $u$ to $v$.
The converse also holds.
\end{lemma}

\begin{proof}
The edge sets with odd-degree vertices exactly $u,v$ form an affine translate of $\Cyc(N)$.
There are $2^{\beta(N)}=2^{s-1}$ such sets.
Every terminal path is one of them.

If equality holds, every such set is a path.
For each terminal path $A$, the graph $N-E(A)$ is a forest.
Otherwise adjoining an edge-disjoint cycle to $A$ gives a join that is not a path.
Since $|E(N)|=2s-2$, this implies $|E(A)|\ge s-1$.
Every terminal path is therefore Hamiltonian.
Fix one such path.
An edge between nonconsecutive vertices would give a shorter terminal path.
Thus every edge joins consecutive vertices of the fixed path.
The degree equations, starting at $u$, force multiplicity two on each consecutive pair.
The converse follows by choosing one edge independently from each doubled pair.
\end{proof}

\begin{lemma}[Expansion]\label{lem:expansion}
Let $H$ be a connected four-regular multigraph on $s$ vertices with two chosen parallel edges between $u,v$.
Delete those edges and replace them by doubled edges $uw$ and $wv$, where $w$ is new.
Call the resulting graph $G$.
Then
\[
 c(G)=c(H)+a+1,\qquad 0\le a\le2^s,
\]
where $a$ counts cycles of $H$ using exactly one chosen edge.
Consequently, the bound in Theorem~\ref{thm:main} for $H$ implies that bound for $G$.
If equality holds for $G$, then $G$ is a doubled cycle.
\end{lemma}

\begin{proof}
Cycles avoiding the chosen pair remain unchanged.
Each cycle using exactly one chosen edge has twice as many replacements.
The chosen digon is replaced by two digons.
This gives the identity.

Let $N$ be obtained by deleting the chosen pair.
If $N$ is disconnected, then $a=0$.
Otherwise $a=2P_{uv}(N)\le2^s$ by Lemma~\ref{lem:paths}.
Equality in the resulting bound for $G$ requires $a=2^s$.
That lemma makes $N$ a doubled path, hence makes both $H$ and $G$ doubled cycles.
\end{proof}

\begin{lemma}[Gluing]\label{lem:gluing}
For $i=1,2$, let $H_i$ be connected and four-regular on $s_i$ vertices.
Choose two parallel edges with distinct ends $u_i,v_i$.
Assume that deleting the chosen pair leaves a connected graph $N_i$.
Delete both chosen pairs and add doubled edges $u_1u_2$ and $v_1v_2$.
Call the resulting graph $G$.
If the bound in Theorem~\ref{thm:main} holds for both $H_i$, it holds for $G$.
Equality for $G$ forces $G$ to be a doubled cycle.
\end{lemma}

\begin{proof}
Put $A_i=2^{s_i}$, $a_i=2P_{u_iv_i}(N_i)$, and $b_i=c(N_i)$.
Lemma~\ref{lem:paths} gives $2\le a_i\le A_i$.
Direct cycle classification gives
\[
 c(H_i)=a_i+b_i+1,\qquad c(G)=a_1a_2+b_1+b_2+2.
\]
A crossing cycle chooses one terminal path in each shore and one edge from each new doubled pair.
All other cycles lie in a shore or are one of the two new digons.
Therefore
\[
 c(G)\le a_1a_2-a_1-a_2+A_1+A_2+s_1+s_2.
\]
The difference from $A_1A_2+s_1+s_2$ is
\[
 (A_1-a_1)(A_2-1)+(A_2-a_2)(a_1-1)\ge0.
\]
Both coefficients are positive, so equality requires $a_i=A_i$ for both $i$.
Lemma~\ref{lem:paths} makes both $N_i$ doubled paths.
Their gluing is a doubled cycle.
The argument permits a chosen pair to belong to a larger parallel class.
\end{proof}

\begin{lemma}[Parity count]\label{lem:parity}
Let $G$ be connected and four-regular.
Suppose its only parallel classes are $k$ doubled pairs, and these pairs form a matching.
Delete their $2k$ edges to obtain $F$, and put $t=\comp(F)$.
If $k-t\ge2$, then
\[
 c(G)\le\frac{27}{32}2^n-1+k<2^n+n.
\]
\end{lemma}

\begin{proof}
Contract the components of $F$ and represent each doubled pair by one edge.
The resulting connected multigraph $J$ has $t$ vertices and $k$ edges; loops are allowed.
Let $Z$ count even edge sets of $G$ using at most one edge from each doubled pair.
A selected edge set of $J$ permits a parity completion inside $F$ exactly when it is even.
Each such selection has $2^{\beta(F)}$ completions and two choices for each selected pair.
Since $\beta(F)=n-2k+t$, a parity filter gives
\[
 Z=2^{n-2k+t}\sum_{S\in\Cyc(J)}2^{|S|}
  =2^{n-2k}\sum_{\sigma\in\{-1,1\}^{V(J)}}
       \prod_{uv\in E(J)}(1+2\sigma_u\sigma_v).
\]
Fix a spanning tree of $J$.
Every non-tree factor has absolute value at most three.
Summing absolute tree products gives $2\cdot4^{t-1}$, by choosing the root sign and then each tree-edge agreement.
Hence
\[
 \frac Z{2^n}\le2^{-2k}\,2\cdot4^{t-1}3^{k-t+1}
       =\frac32\left(\frac34\right)^{k-t}.
\]
Every cycle except the $k$ digons is a nonempty set counted by $Z$.
Thus $c(G)\le Z-1+k$, proving the claim.
\end{proof}

\begin{proposition}[Reduction to small matchings]\label{prop:reduction}
In an induction on the order for Theorem~\ref{thm:main}, every remaining graph is connected, has a Hamilton decomposition,
and has at most two doubled pairs forming a matching.
Its underlying simple graph has minimum degree at least three.
\end{proposition}

\begin{proof}
Every nonempty component of a loopless four-regular graph has at least two vertices.
For a disconnected graph, the bounds on smaller components give a strict inequality because
$\sum_i2^{n_i}<2^{\sum_i n_i}$.
The two-vertex connected graph has four parallel edges and six cycles, as required.
We can assume that $G$ is connected and $n\ge3$.

The $2^{n+1}$ even edge sets form $2^n$ complementary pairs.
A pair containing two cycles must consist of two Hamilton cycles: their degrees sum to four at every vertex.
If no such pair exists, then $c(G)\le2^n$, which is strict for the desired bound.
Thus assume that $G$ has a Hamilton decomposition.
It is four-edge-connected, and every multiplicity is at most two.

Suppose a vertex $w$ has only two distinct neighbours $u,v$.
Each Hamilton cycle uses one edge to each neighbour at $w$, so both multiplicities are two.
For $n=3$, the graph is the doubled triangle.
For $n>3$, no $uv$ edge exists: such an edge would close a triangle in its Hamilton cycle.
Suppress $w$ and replace its four edges by a doubled $uv$ edge.
The resulting smaller graph satisfies the theorem by induction.
Lemma~\ref{lem:expansion} handles the original graph, including equality.
Thus the remaining underlying simple graph has minimum degree at least three.
Every vertex then meets at most one doubled pair.

If there are no doubled pairs, the reduction is complete.
Otherwise define $F,J,k,t$ as in Lemma~\ref{lem:parity}.
Each Hamilton cycle uses exactly one edge from every doubled pair.
Its restriction to a component of $F$ shows that the number of incident matching endpoints is even.
Thus $J$ is Eulerian.
The parity lemma disposes of $k-t\ge2$ with strict inequality.
Assume $k-t\le1$.

If $t>1$, every vertex of $J$ has positive even non-loop degree.
Some vertex has non-loop degree two; otherwise $k\ge2t>t+1$.
Its two non-loop edges form a cut with connected shores.
Indeed, each component on the complementary shore needs a positive even number of these two cut edges.
The corresponding cut of $G$ consists of two doubled pairs with four distinct endpoints.
Close each shore by a new doubled pair between its boundary vertices.
This produces smaller connected loopless four-regular graphs, with connected shores after deletion of the new pair.
Apply induction and Lemma~\ref{lem:gluing}, including its equality conclusion.
We may therefore assume $t=1$.
Now $k-t\le1$ gives $k\le2$.
\end{proof}

\section{The degree-sequence input}

We use Theorem~7 of~\cite{dmnnp}:
for a simple graph $T$ and distinct terminals $s,t$,
\[
 P_{st}(T)\le a(\deg_T(s))\prod_{v\ne s}b(\deg_T(v)).
\]
The functions are nonnegative and the theorem permits disconnected graphs.
Only the following values are needed:
\[
 \rho=(4/3)^{1/4},\quad b(1)=1,\quad b(2)=\rho,\quad
 b(3)=(2\rho)^{1/2},\quad b(4)=2/\rho,\quad b(5)=(24\rho)^{1/4},
\]
\[
 a(1)=\sqrt{2/\rho^5},\quad a(2)=2/\rho^2,\quad a(3)=6/(\rho^4b(5)).
\]
Put
\[
 \kappa=\frac{a(3)b(3)}{b(4)^2}=\frac{3^{19/16}}{2^{17/8}},\qquad
 \delta=\rho^{-2},\qquad
 \chi=\frac2{\sqrt3}+\frac38(\rho/6)^{1/4}.
\]
For a four-regular graph with $k$ matching doubled pairs, put
\[
 B=b(4)^n\rho^k=2^n(3/4)^{(n-k)/4}.
\]
Subdividing one edge from each doubled pair produces a simple graph with degree product $B$.
Subdivision preserves cycles and paths between original terminals.

\begin{lemma}[Edge frequencies]\label{lem:frequency}
Let $T$ be simple with all degrees in $\{2,3,4\}$.
Put $Q=\prod_vb(\deg_T(v))$.
If no edge joins two degree-two vertices, then $c_e(T)\le\delta Q$ for every edge $e$.
Without that restriction, $c_e(T)\le\delta_2Q$, where $\delta_2=\sqrt{2/\rho^9}$.
For a simple four-regular graph, the sharper bound is $c_e(G)\le\kappa B$.
\end{lemma}

\begin{proof}
For endpoint degrees $d,e$, delete the edge and apply the path theorem in either orientation.
Its cycle frequency is at most
\[
 \frac{\min\{a(d-1)b(e-1),a(e-1)b(d-1)\}}{b(d)b(e)}Q.
\]
The following coefficients suffice:
\[
\begin{array}{c|rrrrrr}
(d,e)&(2,2)&(2,3)&(2,4)&(3,3)&(3,4)&(4,4)\\\hline
\text{coefficient}&\delta_2&\rho^{-3}&\rho^{-2}&\rho^{-2}&\kappa/\rho&\kappa.
\end{array}
\]
All coefficients except the first are at most $\delta$, and $\delta<\delta_2$.
\end{proof}

\begin{lemma}[Deleting one vertex]\label{lem:deletion}
Suppose $G$ is four-regular, its doubled pairs form a matching, and all other edges are single.
Write $N=G-u$.
If $k=0$, choose any $u$; then
\[
 c(G)-c(N)\le2\kappa B,\qquad
 \sum_{C\subseteq N}|E(C)|\le(2n-4)\frac\rho6B.
\]
If $k\ge1$, choose $u$ incident with a doubled pair; then
\[
 c(G)-c(N)\le1+\chi B,\qquad
 \sum_{C\subseteq N}|E(C)|\le(2n-4)\frac{\rho^3}{8}B.
\]
In the latter case, $N$ has $k-1$ doubled pairs.
\end{lemma}

\begin{proof}
For $k=0$, each cycle through $u$ uses two incident edges, so their number is at most $2\kappa B$.
The graph $N$ has four vertices of degree three and $n-5$ of degree four.
Its product is $Q=b(3)^4b(4)^{n-5}=(\rho^7/8)B$.
Apply Lemma~\ref{lem:frequency} to its $2n-4$ edges, using $\delta\rho^7/8=\rho/6$.

For $k\ge1$, let $uv$ be the doubled pair, and let $x,y$ be the other neighbours of $u$.
These four vertices are distinct.
In $N$, the degrees at $v,x,y$ are $2,3,3$; all other original degrees are four.
Subdivide one edge from every remaining doubled pair.
The resulting simple graph $T$ has
\[
 Q=b(3)^2b(4)^{n-4}\rho^k=\frac{\rho^5}{8}B.
\]
Its degree-two vertices are $v$ and the subdivision vertices, with no edges between them.
For each original doubled pair, its unsubdivided representative has the same cycle frequency as either original edge.
Thus every original edge of $N$ has frequency at most $\delta Q$.
Their sum gives the length bound because $\delta\rho^5/8=\rho^3/8$.

The cycles through $u$ satisfy
\[
 c(G)-c(N)=1+2P_{vx}(N)+2P_{vy}(N)+P_{xy}(N).
\]
Apply the path theorem to $T$ to bound this expression by
\[
 1+\bigl(4a(2)b(3)^2+a(3)\rho b(3)\bigr)b(4)^{n-4}\rho^{k-1}.
\]
The coefficient of $B$ is
\[
 \frac{4a(2)b(3)^2+a(3)\rho b(3)}{b(4)^4\rho}
 =\rho^2+\frac{3b(3)}{8b(5)}=\chi.
\]
\end{proof}

\begin{lemma}[Short cycles]\label{lem:short}
Let $W$ have maximum degree at most four, with $m$ vertices and $j$ doubled pairs forming a matching.
Assume that all other edges are single.
For $\ell\ge3$,
\[
 c_\ell(W)\le\frac{(6m+2j)3^{\ell-3}}\ell.
\]
\end{lemma}

\begin{proof}
A starting vertex outside the doubled matching gives at most $4\cdot3^{\ell-2}$ directed cycle representations.
There are four first-edge choices, at most three intermediate choices, and at most one closing edge.
At a starting vertex incident with a doubled pair, cycles using two single edges give at most $2\cdot3^{\ell-2}$ representations.
If the first edge is doubled, there are two first-edge choices and at most two second-edge choices.
The closing edge must be single, giving at most $4\cdot3^{\ell-3}$ representations.
Reversal gives the same count when only the closing edge is doubled.
The total at such a vertex is at most $14\cdot3^{\ell-3}$.
Sum over the $m-2j$ and $2j$ starting vertices and divide by $2\ell$.
\end{proof}

\section{All reduced graphs of order at least eleven}

\begin{proposition}\label{prop:large}
If $G$ is four-regular, $n\ge11$, and its only parallel classes are at most two matching doubled pairs, then $c(G)<2^n$.
\end{proposition}

\begin{proof}
First consider $k=0,n\ge12$, or $k\in\{1,2\},n\ge11$.
Put $g=\lfloor(n+3)/2\rfloor$, $j=\max\{k-1,0\}$, and
\[
 F_g=\sum_{\ell=3}^{g-1}(1-\ell/g)\frac{3^{\ell-3}}\ell.
\]
Apply the length cutoff to $N=G-u$ from Lemma~\ref{lem:deletion}.
Every cycle contributes its length divided by $g$, supplemented by $1-\ell/g$ when its length is less than $g$.
Lemma~\ref{lem:short} yields
\begin{equation}\label{eq:analytic-bound}
 \frac{c(G)}{2^n}\le A_{n,k}+T_{n,k}(3/4)^{(n-k)/4},
\end{equation}
where
\[
 A_{n,k}=2^{-n}\left({\bf1}_{k\ge1}+(1-2/g)j+(6(n-1)+2j)F_g\right)
\]
and
\[
 T_{n,k}=\begin{cases}
 2\kappa+(2n-4)\rho/(6g),&k=0,\\
 \chi+(2n-4)\rho^3/(8g),&k\ge1.
 \end{cases}
\]
Use these rational upper bounds:
\begin{equation}\label{eq:rational-constants}
 \kappa<\frac{845071}{10^6},\quad
 \rho<\frac{1074570}{10^6},\quad
 \rho^3<\frac{1240807}{10^6},\quad
 \chi<\frac{1398653}{10^6}.
\end{equation}
Write $\overline T_{n,k}$ for the resulting rational coefficient.
For the six base pairs below, exact arithmetic verifies
\[
 A_{n,k}<1,\qquad
 \overline T_{n,k}^{\,4}(3/4)^{n-k}<(1-A_{n,k})^4.
\]
The final column gives a rounded-up upper bound from~\eqref{eq:analytic-bound}.
\[
\begin{array}{c|c|c|c}
 n&k&g&\text{upper ratio}\\\hline
 12&0&7&.956\\
 13&0&8&.891\\
 11&1&7&.926\\
 12&1&7&.863\\
 11&2&7&.993\\
 12&2&7&.926
\end{array}
\]

We show that the bound decreases under $n\mapsto n+2$ with $k$ fixed.
Then $g\mapsto g+1$.
For $g\ge6$, we have $F_{g+1}\le3F_g$.
To verify this, shift the index of $3F_g$ by one and compare terms with indices $4,\ldots,g$.
Each term dominates the corresponding term of $F_{g+1}$.
The excess at index $g$ alone is
\[
 \frac{2\cdot3^{g-3}}{g(g^2-1)}
 \ge\frac{g-2}{3(g+1)},
\]
which covers its missing index-three term.
The inequality is equivalent to $2\cdot3^{g-2}\ge g(g-1)(g-2)$.
It holds at $g=6$ and persists because the right-hand growth ratio is $(g+1)/(g-2)<3$.

The coefficient of $F_g$ increases by at most $(n+1)/(n-1)$.
The other nonnegative terms in the numerator of $A_{n,k}$ increase by a factor at most three.
Consequently,
\[
 A_{n+2,k}\le\frac{3(n+1)}{4(n-1)}A_{n,k}<A_{n,k}\qquad(n\ge11).
\]
For the long term, write $\overline T=a+b(n-2)/g$, where $a>1$ and $0<b<1$.
Then
\[
 \overline T_{n+2,k}-\overline T_{n,k}
 =b\frac{2g-n+2}{g(g+1)}\le\frac5{56}<\frac1{10}.
\]
Thus $\overline T_{n+2,k}<(11/10)\overline T_{n,k}$.
The exponential factor decreases by $\sqrt3/2$, and $(11/10)^2(3/4)=363/400<1$.
Both terms therefore decrease, proving all these orders from the six base cases.

It remains to treat a simple graph on eleven vertices.
Choose adjacent $u,v$, and put $N=G-u$, $H=G-\{u,v\}$.
Cycles through $u$ contribute at most $2\kappa B$.
In $N$, the vertex $v$ has degree three and $\delta Q_N=(\rho/6)B$.
Thus cycles through $v$ in $N$ contribute at most $(\rho/4)B$.

If $u,v$ have $r$ common neighbours, then $0\le r\le3$.
The degrees in $H$ are two at $r$ vertices, three at $6-2r$ vertices, and four at $3+r$ vertices.
Consequently,
\[
 \frac{Q_H}{B}=\frac{\rho^r b(3)^{6-2r}}{b(4)^{8-r}}
 =\frac{\rho^{11-r}}{32}\le\frac{\rho^3}{18}.
\]
There are fifteen edges in $H$.
Lemma~\ref{lem:frequency} gives their total cycle frequency at most $(5/6)\delta_2\rho^3B$.
Apply cutoff six in $H$, using Lemma~\ref{lem:short} with $m=9,j=0$.
This gives
\[
 c(G)\le\left(2\kappa+\frac\rho4+\frac5{36}\delta_2\rho^3\right)B
       +\sum_{\ell=3}^{5}(1-\ell/6)\frac{18\cdot3^{\ell-2}}\ell.
\]
Use~\eqref{eq:rational-constants} and $\delta_2<1023197/10^6$.
An exact fourth-power comparison gives $c(G)/2^{11}<.987$.
This completes the proof.
\end{proof}

The constant comparisons above use only rational powers.
For clarity, $\kappa^{16}=3^{19}/2^{34}$ and $\delta_2^8=16(3/4)^9$.
For $\chi$, use $1/\sqrt3<577351/10^6$ and
$\frac38(\rho/6)^{1/4}<243951/10^6$, verified by squares and sixteenth powers.
The supplied certificate script checks all seven base cases and the induction constants exactly.
It also records seventeen additional finite checks that are not needed for the induction.

\section{Finite verification and completion of the proof}

The remaining graphs have minimum underlying degree three, at most two matching doubled pairs, and $n\le10$.
For each $k\in\{0,1,2\}$, remove one edge from every doubled pair.
The resulting connected simple graph has $2n-k$ edges, degree three at exactly $2k$ vertices, and degree four elsewhere.
Conversely, every desired multigraph arises by doubling a perfect matching of existing edges among these degree-three vertices.
For $k=1$ there is at most one choice; for $k=2$ there are at most three.

We generated all such underlying graphs with
\begin{verbatim}
geng -c -d3 -D4 n (2*n-k)
\end{verbatim}
where the last argument is evaluated as an integer.
For each graph, the checker tries every permitted matching and counts all cycles exactly.
One implementation uses vertex-simple cyclic orders and multiplies edge multiplicities.
The other constructs a fundamental binary cycle basis and retains precisely the connected nonzero degree-two edge sets.
These methods distinguish the parallel edges in the same way.

The following table records each nonempty class.
Entries are the maximum cycle count and, in parentheses, the number of matching instances.
Different matchings may produce isomorphic multigraphs; duplicates do not affect completeness.
\[
\begin{array}{c|rrr}
 n&k=0&k=1&k=2\\\hline
 4&-&-&19\ (3)\\
 5&37\ (1)&-&32\ (2)\\
 6&63\ (1)&60\ (1)&52\ (2)\\
 7&108\ (2)&102\ (3)&97\ (5)\\
 8&204\ (6)&177\ (10)&165\ (22)\\
 9&333\ (16)&325\ (41)&286\ (90)\\
 10&593\ (59)&534\ (200)&526\ (461)
\end{array}
\]
All $925$ instances satisfy $c(G)\le2^n+n$.
The only equality instance is $K_5$.
The generator considered $3234$ underlying graphs, including those without a permitted matching.
The two counting methods also agree on all $5012$ matching instances through order eleven.
Only the $925$ instances through order ten are needed for the proof.

\begin{proof}[Proof of Theorem~\ref{thm:main}]
Use induction on $n$.
The empty case is strict; no one-vertex loopless four-regular graph exists.
The two-vertex graph has four parallel edges and six cycles.
Proposition~\ref{prop:reduction} handles disconnected graphs, graphs without a Hamilton decomposition,
doubled-edge expansions, and the parity and gluing cases.
Each reduction either gives a strict bound or identifies a doubled cycle when equality holds.

The remaining graphs have at most two matching doubled pairs and minimum underlying degree three.
Proposition~\ref{prop:large} gives a strict bound for $n\ge11$.
The finite verification handles all smaller remaining graphs, with equality only for $K_5$.
Finally, a doubled cycle has $2^n+n$ cycles, and
\[
 c(K_5)=\binom53+3\binom54+\frac{4!}{2}=10+15+12=37=2^5+5.
\]
This proves both the inequality and its equality classification.
\end{proof}


\section{The graphs above the one-half bound}

Let $D$ be $K_4$ with two opposite edges doubled.
It has $19$ cycles: two digons, eight triangles, and nine Hamilton cycles.
For either doubled pair, exactly twelve cycles use one edge from that pair.
Indeed, deleting the pair leaves six paths between its endpoints.

\begin{theorem}\label{thm:above-half}
For a loopless four-regular multigraph $G$ on $n$ vertices,
\[
 c(G)>2^n
\]
holds if and only if $G$ is a doubled cycle, $K_5$, or $D$.
\end{theorem}

\begin{proof}
These three families have the stated cycle counts.
For the converse, use induction on the order and the reductions from Proposition~\ref{prop:reduction}.
A disconnected graph has $c(G)<2^n$.
Indeed, Theorem~\ref{thm:main} applies to each component, and
$2^a+2^b+a+b<2^{a+b}$ for integers $a,b\ge2$.
A graph without a Hamilton decomposition also has $c(G)\le2^n$.

Consider an expansion of a smaller graph $H$, with $A=2^{|V(H)|}$ and the notation of Lemma~\ref{lem:expansion}.
If $H$ is a doubled cycle, its expansion is also a doubled cycle.
If $H=D$, then $c(G)=19+12+1=32=2^5$.
The graph $K_5$ has no doubled pair and cannot occur here.
Every other $H$ has $c(H)\le A$ by induction.
By Lemma~\ref{lem:paths}, its even integer $a$ is at most $A-2$.
Therefore $c(G)\le2A-1<2^{|V(G)|}$.

The parity case gives $c(G)\le(27/32)2^n-1+k<2^n$.
Here $k\le n/2$, and the displayed strict inequality follows from $2^n\ge4n$ for $n\ge4$.
It remains to check the gluing reduction.
Keep $A_i,a_i$ from Lemma~\ref{lem:gluing} and write $c_i=c(H_i)$.
Then
\[
 c(G)=a_1a_2+c_1+c_2-a_1-a_2.
\]
Call a shore ordinary if $c_i\le A_i$.
An ordinary shore is not a doubled cycle, so $a_i\le A_i-2$.
The other possibilities are doubled cycles and $D$, because a shore must have a chosen parallel pair.

Two ordinary shores give $c(G)\le A_1A_2$ by the same monotonicity used in Lemma~\ref{lem:gluing}.
A doubled cycle on $s_1$ vertices and an ordinary shore give
\[
 c(G)\le(A_1-1)(A_2-2)+A_2+s_1
       =A_1A_2-2A_1+s_1+2<A_1A_2.
\]
A $D$ shore and an ordinary shore give $c(G)\le12A_2+7<16A_2$.
A $D$ shore and a doubled cycle give $c(G)=12A_2+s_2+7<16A_2$.
Two $D$ shores give $c(G)=158<256$.
Two doubled-cycle shores glue to a doubled cycle.
This completes the gluing cases.

The remaining reduced graphs have $k\le2$.
Proposition~\ref{prop:large} puts every order at least eleven strictly below $2^n$.
The finite table shows that only $K_5$ and $D$ exceed $2^n$ at the remaining orders.
This proves the classification.
\end{proof}

\section{Consequences for Eulerian graphs}

\begin{corollary}\label{cor:eulerian}
Let $G$ be a connected loopless Eulerian multigraph of cycle rank $r\ge1$.
Then
\[
 c(G)\le2^{r-1}+r-1.
\]
For $r\ge3$, equality holds precisely for subdivisions of doubled cycles and of $K_5$.
The $K_5$ case has rank six.
At $r=1$, the graph is a single cycle and equality holds.
At $r=2$, the maximum is two cycles, so the displayed bound is strict.
Moreover, for $r\ge3$,
\[
 c(G)>2^{r-1}
\]
holds precisely for subdivisions of doubled cycles, $K_5$, or $D$.
\end{corollary}

\begin{proof}
Since $G$ is Eulerian, complementation pairs its $2^r$ even edge sets.
If $c(G)>2^{r-1}$, one pair consists of two cycles $C_1,C_2$ partitioning $E(G)$.
Every vertex then has degree two or four.
The degree-four vertices are exactly the common vertices of $C_1,C_2$.
If their number is $s$, then $|E(G)|=|V(G)|+s$ and $r=s+1$.

Connectedness implies $s\ge1$.
If $s=1$, the graph consists of two cycles meeting only at their common vertex.
It has exactly two cycles, contradicting $c(G)>2^{r-1}=2$.
Hence $s\ge2$ whenever the one-half bound is exceeded.
Suppress every degree-two vertex.
Each $C_i$ meets all $s$ common vertices exactly once.
Its segments therefore have distinct endpoints, so suppression creates no loops.
The resulting graph $H$ is loopless and four-regular on $s=r-1$ vertices.
Suppression preserves all cycles.

Theorem~\ref{thm:main} gives
$c(G)=c(H)\le2^{r-1}+r-1$.
If the one-half bound was not exceeded, the same inequality already holds.
For $r\ge3$, equality in the new bound exceeds the one-half bound, so its classification follows from Theorem~\ref{thm:main}.
The classification above the one-half bound follows from Theorem~\ref{thm:above-half}.
Conversely, subdivision preserves cycle rank and cycle count, so every listed subdivision has the stated property.

At rank one, a connected Eulerian graph is a cycle.
At rank two, more than two cycles would trigger the same complementary-pair argument with $s=1$, a contradiction.
Two cycles meeting at one vertex attain the value two.
\end{proof}

\section{Reproducibility and proof status}

The mathematical proof uses the published path theorem~\cite{dmnnp} and a finite computer verification.
The computation is not a Lean certificate.
The relevant files beside this manuscript are:
\begin{itemize}
\item \nolinkurl{four_regular_vertex_deletion_certificates.py} verifies the exact rational analytic inequalities.
Its output is \nolinkurl{four-regular-vertex-deletion-results.json}.
\item \nolinkurl{check_four_regular_finite.cpp} counts cycles by vertex-simple orders.
\item \nolinkurl{verify_four_regular_theorem.py} generates the graphs and independently counts their binary cycle spaces.
\item \nolinkurl{four-regular-theorem-certificate.json} records the complete finite verification.
\end{itemize}
The latter record includes generator commands and implementation hashes.
The nauty source archive is
\url{https://pallini.di.uniroma1.it/nauty2_8_9.tar.gz}, with SHA-256
\begin{verbatim}
c97ab42bf48796a86a598bce3e9269047ca2b32c14fc23e07208a244fe52c4ee
\end{verbatim}
Build its \texttt{geng} target, then compile the checker with a C++17 compiler.
The runner accepts explicit paths through \texttt{--geng} and \texttt{--checker}.
The default maximum order ten reproduces the exact small-core table.
The optional \texttt{--max-n 11} also checks the next order independently.
This runner uses no analytic graph-count certificate; it compares both exact cycle counts on every instance.

\begin{thebibliography}{9}
\bibitem{aacst}
A. Alahmadi, R. E. L. Aldred, R. dela Cruz, P. Sol\'e, and C. Thomassen,
\emph{The maximum number of minimal codewords in long codes},
arXiv:1201.6043, Section~6, Open Problem~2.
\url{https://arxiv.org/abs/1201.6043}.
\bibitem{dmnnp}
Z. Dvo\v r\'ak, N. Morrison, J. A. Noel, S. Norin, and L. Postle,
\emph{Bounding the number of cycles in a graph in terms of its degree sequence},
arXiv:1907.12091, Theorem~7.
\url{https://arxiv.org/abs/1907.12091}.
\bibitem{nauty}
B. D. McKay and A. Piperno,
\emph{nauty and Traces}, version~2.8.9, including the \texttt{geng} graph generator.
\url{https://pallini.di.uniroma1.it/}.
\end{thebibliography}
\end{document}
